\section{本書で見たつながり}\label{節：本書で見たつながり}

連立漸化式から出発して,隣接多項間漸化式,複素平面上の回転,一次分数型へ進んだ.
最後には,一定角の回転から角度を倍にする更新へ範囲を広げた.
一方で,\,式の見た目が似ていても更新の仕組みが異なる場合もある.

\begin{itemize}
 \item $z_{n+1}=\lambda z_n$は,\,一定角の回転と一定倍率の拡大・縮小である.
 \item $z_{n+1}=z_n^d$は,\,偏角を$d$倍し,\,絶対値も$d$乗する. 単位円上でも一定角の回転とは異なる.
 \item 実部や傾きだけを記録すると,\,異なる点が同じ値になることがある. 何の情報を残し,\,何を区別しなくなったかを確認する.
\end{itemize}

複素平面への置き換えが,\,常に全体の座標変換になるわけではない.
例えば余弦は上下二点を同じ横座標に送る.
それでも,\,元の更新と整合する点の動きを一つ作れば,\,数列の式を導ける.
一対一の座標変換と,\,一部の情報を読み取る操作を区別して使おう.

\begin{bookpoint}
漸化式の形は,\,動きそのものだけでなく,\,\textbf{その動きの何を数として記録したか}によって決まる.
不動点からのずれ,\,固有方向の成分,\,成分比,\,実部,\,傾き.
追う量を選び直すと,\,別々に見えた解法のつながりが見えてくる.
\end{bookpoint}
